\Needspace{15\baselineskip}
\Article{25}

\textsc{Satz.}\quad \textit{Es sei $n$ das Produkt zweier teilerfremder positiver ganzer Zahlen $a$, $b$,} \textit{und es werde gesetzt}

\[
\begin{aligned}
P&=1+ r^{aa}+r^{4aa}+r^{9aa}+\text{usw.}+r^{(b-1)^2aa}\\
Q&=1+r^{bb}+r^{4bb}+r^{9bb}+\text{usw.}+r^{(a-1)^2bb}
\end{aligned}
\]

\textit{Dann behaupte ich, dass $W=PQ$ ist.}

\textit{Beweis.}\quad Es bezeichne $\alpha$ unbestimmt die Zahlen $0,1,2,3\dotsfour a-1$, $\beta$ die Zahlen $0,1,2,3\dotsfour b-1$ und $\nu$ die Zahlen $0,1,2,3\dotsfour n-1$.\quad Dann gilt offenbar

\[
\begin{aligned}
P=\Sigma r^{aa\beta\beta},\quad Q=\Sigma r^{bb\alpha\alpha},\quad W=\Sigma r^{\nu\nu}
\end{aligned}
\]

Somit ist $PQ=\Sigma r^{aa\beta\beta+bb\alpha\alpha}$, wenn für $\alpha$ und $\beta$ sämtliche Werte in allen möglichen Kombinationen eingesetzt werden. Da ferner $2ab\alpha\beta=2\alpha\beta n$ ist, ergibt sich $PQ=\Sigma r^{(a\beta+b\alpha)^2}$. Man sieht aber ohne Mühe, dass die einzelnen Werte von $a\beta+b\alpha$ voneinander verschieden und jeweils einem Wert von $\nu$ gleich sind.\quad Daher ist $PQ=\Sigma r^{\nu\nu}=W$.

Außerdem ist zu beachten, dass $r^{aa}$ eine eigentliche Wurzel der Gleichung $x^b-1=0$ und $r^{bb}$ eine eigentliche Wurzel der Gleichung $x^a-1=0$ ist.

\SourceLeaf{47}

\Needspace{10\baselineskip}
\Article{26}

Es sei nun $n$ das Produkt dreier paarweise teilerfremder Zahlen $a,b,c$. Wird $bc=b'$ gesetzt, so sind offenbar auch $a$ und $b'$ teilerfremd; daher ist $W$ das Produkt der beiden Faktoren

\[
\begin{aligned}
&1+r^{aa}+r^{4aa}+r^{9aa}+\text{usw.}+r^{(b'-1)^2aa}\\
&1+r^{b'b'}+r^{4b'b'}+r^{9b'b'}+\text{usw.}+r^{(a-1)^2b'b'}
\end{aligned}
\]

Da aber $r^{aa}$ eine eigentliche Wurzel der Gleichung $x^{bc}-1=0$ ist, ist der erste Faktor selbst das Produkt aus

\[
\begin{aligned}
&1+\rho^{bb}+\rho^{4bb}+\rho^{9bb}+\text{usw.}+\rho^{(c-1)^2bb}\\
&1+\rho^{cc}+\rho^{4cc}+\rho^{9cc}+\text{usw.}+\rho^{(b-1)^2cc}
\end{aligned}
\]

wenn $r^{aa}=\rho$ gesetzt wird.\quad Daraus ist ersichtlich, dass $W$ das Produkt der drei Faktoren

\[
\begin{aligned}
&1+r^{bbcc}+r^{4bbcc}+r^{9bbcc}+\text{usw.}+r^{(a-1)^2bbcc}\\
&1+r^{aacc}+r^{4aacc}+r^{9aacc}+\text{usw.}+r^{(b-1)^2aacc}\\
&1+r^{aabb}+r^{4aabb}+r^{9aabb}+\text{usw.}+r^{(c-1)^2aabb}
\end{aligned}
\]

ist, wobei $r^{bbcc}$, $r^{aacc}$, $r^{aabb}$ jeweils eigentliche Wurzeln der Gleichungen $x^a-1=0$, $x^b-1=0$, $x^c-1=0$ sind.

\Needspace{12\baselineskip}
\Article{27}

Daraus folgt leicht allgemein: Ist $n$ das Produkt beliebig vieler paarweise teilerfremder Faktoren $a,b,c$ usw., so ist $W$ das Produkt ebenso vieler Faktoren, nämlich

\[
\begin{aligned}
&1+r^{\frac{nn}{aa}}+r^{\frac{4nn}{aa}}+r^{\frac{9nn}{aa}}+\text{usw.}+r^{\frac{(a-1)^2nn}{aa}}\\
&1+r^{\frac{nn}{bb}}+r^{\frac{4nn}{bb}}+r^{\frac{9nn}{bb}}+\text{usw.}+r^{\frac{(b-1)^2nn}{bb}}\\
&1+r^{\frac{nn}{cc}}+r^{\frac{4nn}{cc}}+r^{\frac{9nn}{cc}}+\text{usw.}+r^{\frac{(c-1)^2nn}{cc}},\;\text{usw.}
\end{aligned}
\]

wobei $r^{\frac{nn}{aa}}$, $r^{\frac{nn}{bb}}$, $r^{\frac{nn}{cc}}$ usw. eigentliche Wurzeln der Gleichungen $x^a-1=0$, $x^b-1=0$. $x^c-1=0$ usw. sind.

\Needspace{14\baselineskip}
\Article{28}

Aus diesen Grundsätzen ergibt sich nun unmittelbar der Weg zur vollständigen Bestimmung von $W$ für einen beliebigen Wert von $n$.\quad Zerlegen wir nämlich $n$ in Faktoren
\SourceLeaf{48}
$a,b,c$ usw., die entweder verschiedene Primzahlen oder Potenzen verschiedener Primzahlen sind, und setzen $r^{\frac{nn}{aa}}=A$, $r^{\frac{nn}{bb}}=B$, $r^{\frac{nn}{cc}}=C$ usw., so sind $A,B,C$ usw. eigentliche Wurzeln der Gleichungen $x^a-1=0$, $x^b-1=0$, $x^c-1=0$ usw. und $W$ das Produkt der Faktoren

\[
\begin{aligned}
&1+A+A^4+A^9+\text{usw.}+A^{(a-1)^2}\\
&1+B+B^4+B^9+\text{usw.}+B^{(b-1)^2}\\
&1+C+C^4+C^9+\text{usw.}+C^{(c-1)^2}
\end{aligned}
\]

Diese einzelnen Faktoren lassen sich aber anhand unserer Ausführungen in den Art. $20,21,23$ bestimmen, wodurch auch der Wert des Produkts bekannt wird.\quad Es wird nützlich sein, die Regeln zur Bestimmung dieser Faktoren hier in einer Übersicht zusammenzustellen.\quad Da die Wurzel $A$ $=\frac{kn}{a}\pdot\frac{360^0}{a}$ ist, wird die Summe $1+A+A^4+A^9+\text{usw.}+A^{(a-1)^2}$, die wir mit $L$ bezeichnen, ebenso durch die Zahl $\frac{kn}{a}$ bestimmt, wie in unserer allgemeinen Untersuchung $W$ durch $k$ bestimmt wird.\quad Nun sind zwölf Fälle zu unterscheiden.

I.\quad Ist $a$ eine Primzahl der Form $4\mu+1$, etwa $=p$, oder eine Potenz einer solchen Primzahl mit ungeradem Exponenten und zugleich $\frac{kn}{a}$ ein quadratischer Rest von $p$, so ist $L=+\surd a$.

II.\quad Bleiben die übrigen Bedingungen unverändert und ist $\frac{kn}{a}$ ein quadratischer Nichtrest von $p$, so ist $L=-\surd a$.

III.\quad Ist $a$ eine Primzahl der Form $4\mu+3$, etwa $=p$, oder eine Potenz einer solchen Primzahl mit ungeradem Exponenten und zugleich $\frac{kn}{a}$ ein quadratischer Rest von $p$, so ist $L=+i\surd a$.

IV.\quad Sind die übrigen Bedingungen wie in III und ist $\frac{kn}{a}$ ein quadratischer Nichtrest von $p$, so ist $L=-i\surd a$.

V.\quad Ist $a$ das Quadrat oder eine höhere Potenz einer (ungeraden) Primzahl mit geradem Exponenten, so ist $L=+\surd a$.

VI.\quad Ist $a=2$, so ist $L=0$.

VII.\quad Ist $a=4$ oder eine höhere Zweierpotenz mit geradem Exponenten und hat zugleich $\frac{kn}{a}$ die Form $4\mu+1$, so ist $L=(1+i)\surd a$.

VIII.\quad Sind die übrigen Bedingungen wie in VII und hat $\frac{kn}{a}$ die Form $4\mu+3$, so ist $L=(1-i)\surd a$.

IX.\quad Ist $a=8$ oder eine höhere Zweierpotenz mit ungeradem Exponenten und hat zugleich $\frac{kn}{a}$ die Form $8\mu+1$, so ist $L=(1+i)\surd a$.

\SourceLeaf{49}

X.\quad Sind die übrigen Bedingungen wie in IX und hat $\frac{kn}{a}$ die Form $8\mu+3$, so ist $L=(-1+i)\surd a$.

XI.\quad Bleiben die übrigen Bedingungen unverändert und hat $\frac{kn}{a}$ die Form $8\mu+5$, so ist $L=(-1-i)\surd a$.

XII.\quad Bleiben die übrigen Bedingungen unverändert und hat $\frac{kn}{a}$ die Form $8\mu+7$, so ist $L=(1-i)\surd a$.

\Needspace{17\baselineskip}
\Article{29}

Es sei beispielsweise $n=2520=8\pdot9\pdot5\pdot7$ und $k=13$.\quad Dann erhalten wir

{\small
\[
\begin{aligned}
&\text{für }a=8,\;\text{nach Fall XII,}\quad L=(1-i)\surd8\\
&\text{für den Faktor }9,\;\text{nach Fall V, zugehörige Summe }=\surd9\\
&\text{für den Faktor }5,\;\text{nach Fall II, zugehörige Summe }=-\surd5\\
&\text{für den Faktor }7,\;\text{nach Fall III, zugehörige Summe }=+i\surd7\mathord{.}
\end{aligned}
\]
}

Somit ist $W=(1-i)\pdot(-i)\pdot\surd2520=(-1-i)\surd2520$.

Für denselben Wert von $n$ sei $k=1$; dann entsprechen den Faktoren die folgenden Summen:

\[
\begin{aligned}
&\text{zum Faktor }8\;\text{die Summe }(-1+i)\surd8\\
&\text{zum Faktor }9\;\text{die Summe }\surd9\\
&\text{zum Faktor }5\;\text{die Summe }\surd5\\
&\text{zum Faktor }7\;\text{die Summe }-i\surd7
\end{aligned}
\]

Daraus ergibt sich das Produkt $W=(1+i)\surd2520$.
